%% KSP 7.0 Selection — Theory Assignment
%% Krittika, IIT Bombay | May 2026
%%
%% Compilation: pdflatex (twice for cross-references)
%% Required packages: amsmath, amssymb, geometry, graphicx,
%%   hyperref, booktabs, xcolor, enumitem, physics, cleveref

\documentclass[11pt, a4paper]{article}

\usepackage[margin=2.4cm, top=2.8cm, bottom=2.8cm]{geometry}
\usepackage{amsmath, amssymb, amsfonts}
\usepackage{physics}
\usepackage{graphicx}
\usepackage{booktabs}
\usepackage{xcolor}
\usepackage{hyperref}
\usepackage{enumitem}
\usepackage{microtype}
\usepackage{cleveref}
\usepackage{caption}

\hypersetup{colorlinks=true, linkcolor=blue!60!black, citecolor=green!50!black,
            urlcolor=cyan!70!black}

\definecolor{noteblue}{RGB}{220,235,250}
\definecolor{noteborder}{RGB}{60,110,180}

\newcommand{\thetaE}{\theta_{\!E}}
\newcommand{\unit}[1]{\,\mathrm{#1}}
% \dd already defined by physics package
\newcommand{\half}{\tfrac{1}{2}}

% Compact note box — used sparingly
\newenvironment{notebox}{%
  \begin{trivlist}\item\relax
  \begin{lrbox}{\@tempboxa}
    \begin{minipage}{0.94\linewidth}
    \color{noteborder}\small}{%
    \end{minipage}
  \end{lrbox}
  \colorbox{noteblue}{\usebox{\@tempboxa}}
  \end{trivlist}}

\setlength{\parindent}{0pt}
\setlength{\parskip}{6pt}

% ─────────────────────────────────────────────────────────────────────────────
\begin{document}

\begin{center}
  {\Large \textbf{KSP 7.0 Selection Assignment — Theory Solutions}}\\[4pt]
  {\large Krittika, IIT Bombay \quad|\quad May 2026}
\end{center}

\vspace{4pt}
\hrule
\vspace{10pt}

\tableofcontents
\newpage

% =============================================================================
\section{Question 1 — Orbital Mechanics: Kepler's Laws and Vis-Viva}

\subsection*{Setup}

Consider a spacecraft of mass $m$ in an elliptical orbit around the Sun (mass $M_\odot$),
with semi-major axis $a$, eccentricity $e$, perihelion distance $r_p = a(1-e)$, and
aphelion distance $r_a = a(1+e)$.

\subsection*{(a) Orbital energy and vis-viva}

The specific mechanical energy of any Keplerian orbit is
\begin{equation}
  \mathcal{E} = -\frac{GM_\odot}{2a},
\end{equation}
which follows directly from the definition $\mathcal{E} = v^2/2 - GM_\odot/r$ with
the constraint that $\mathcal{E}$ is constant.  Equating this at an arbitrary point
$(r, v)$ on the orbit gives the vis-viva equation:
\begin{equation}
  \boxed{v^2 = GM_\odot\!\left(\frac{2}{r} - \frac{1}{a}\right).}
  \label{eq:visviva}
\end{equation}
At perihelion the velocity is maximised ($v_p$) and at aphelion minimised ($v_a$):
\begin{equation}
  v_p = \sqrt{\frac{GM_\odot(1+e)}{a(1-e)}}, \qquad
  v_a = \sqrt{\frac{GM_\odot(1-e)}{a(1+e)}}.
\end{equation}

\subsection*{(b) Kepler's Third Law}

From Eq.~\eqref{eq:visviva} and the geometry of the ellipse the orbital period is
\begin{equation}
  \boxed{T^2 = \frac{4\pi^2}{GM_\odot}\,a^3.}
  \label{eq:k3}
\end{equation}
\textit{Dimensional check:}
$[GM_\odot] = \unit{m^3\,s^{-2}}$, $[a^3] = \unit{m^3}$, so
$[a^3/(GM_\odot)] = \unit{s^2}$, confirming $[T^2] = \unit{s^2}$ \checkmark.

Using $GM_\odot = 1.327 \times 10^{20}\unit{m^3\,s^{-2}}$ and $a_\oplus = 1\unit{AU}
= 1.496 \times 10^{11}\unit{m}$:
\[
  T_\oplus = 2\pi\sqrt{\frac{a_\oplus^3}{GM_\odot}}
           = 2\pi\sqrt{\frac{(1.496\times10^{11})^3}{1.327\times10^{20}}}
           \approx 3.156\times10^7\unit{s} \approx 365.25\unit{days.}
\]

\subsection*{(c) Circular-orbit limit}

For $e \to 0$: $r_p \to r_a \to a$, $v_p \to v_a \to v_c = \sqrt{GM_\odot/a}$
and $T \to 2\pi a/v_c$.  The orbit is a circle with the Sun at the centre and
all the usual Keplerian relations become exact.

% =============================================================================
\section{Question 2 — Hohmann Transfer}

\subsection*{Setup and $\Delta v$ budget}

A Hohmann transfer from a circular orbit of radius $R_1$ to one of radius $R_2 > R_1$
uses two impulsive burns on a half-ellipse with semi-major axis $a_T = (R_1+R_2)/2$.

Initial circular speed: $v_1 = \sqrt{GM_\odot/R_1}$.

Perihelion speed of transfer ellipse:
\[
  v_p = \sqrt{\frac{2GM_\odot R_2}{R_1(R_1+R_2)}}.
\]
First burn:
\begin{equation}
  \Delta v_1 = v_p - v_1 = \sqrt{\frac{GM_\odot}{R_1}}\!\left(\sqrt{\frac{2R_2}{R_1+R_2}}-1\right).
\end{equation}

Aphelion speed of transfer ellipse:
\[
  v_a = \sqrt{\frac{2GM_\odot R_1}{R_2(R_1+R_2)}}.
\]
Second burn (to circularise at $R_2$):
\begin{equation}
  \Delta v_2 = v_2 - v_a = \sqrt{\frac{GM_\odot}{R_2}}\!\left(1-\sqrt{\frac{2R_1}{R_1+R_2}}\right).
\end{equation}

Transfer time $= $ half the period of the transfer ellipse:
\begin{equation}
  t_{\rm transfer} = \pi\sqrt{\frac{a_T^3}{GM_\odot}}
                   = \frac{\pi}{\sqrt{GM_\odot}}\left(\frac{R_1+R_2}{2}\right)^{3/2}.
\end{equation}

\subsection*{Earth–Jupiter Hohmann (numerical)}

$R_1 = 1\unit{AU}$, $R_2 = 5.2\unit{AU}$ (semi-major axis of Jupiter's orbit).
\begin{align*}
  a_T      &= 3.1\unit{AU} = 4.638\times10^{11}\unit{m},\\
  \Delta v_1 &= 8.79\unit{km/s},\\
  \Delta v_2 &= 5.64\unit{km/s},\\
  t_1      &= 2.73\unit{yr}.
\end{align*}
These are independently verified by numerical integration in the coding assignment
(Section~``Validation 2'').

\subsection*{Phase angle requirement}

For the spacecraft to arrive at $R_2$ when Jupiter is also there, Jupiter must be
positioned at an angle $\theta_J$ ahead of Earth at launch:
\[
  \theta_J = \pi - \omega_J\,t_1
           = \pi - \frac{2\pi}{T_J}\cdot t_1.
\]
With $T_J = 11.86\unit{yr}$: $\theta_J = \pi - (0.531\unit{rad}) \approx 2.61\unit{rad}
\approx 149.7^\circ$.  This is the required Jupiter–Earth angular separation at launch.

\subsection*{Connection to Q3 flyby geometry}

The spacecraft exits the Hohmann arc at Jupiter's orbital radius with a heliocentric
speed $v_a$ (aphelion speed of the transfer ellipse).  This speed and direction serve
as the \emph{initial} conditions for the Jupiter-frame encounter analysed in Q3.
Specifically, the angle $\theta_i = 60^\circ$ quoted in Q5 is the angle between
the incoming heliocentric velocity and Jupiter's orbital velocity at the moment of
closest approach.  The flyby scattering angle $\Delta$ computed in Q3 then determines
the post-flyby heliocentric direction, which sets the trajectory to Neptune.

% =============================================================================
\section{Question 3 — Hyperbolic Flyby and Gravity Assist}

\subsection*{Scattering angle}

In Jupiter's rest frame, the spacecraft approaches on a hyperbolic trajectory with
speed $u_\infty$ (the asymptotic speed at infinity) and impact parameter $b$.
For a gravitational interaction with $GM_J$, the geometry of the hyperbola gives the
deflection angle:
\begin{equation}
  \boxed{\Delta = 2\arctan\!\left(\frac{GM_J}{b\,u_\infty^2}\right).}
  \label{eq:delta}
\end{equation}
This follows from the relation between the asymptotic directions of the two branches
of the hyperbola: the semi-transverse axis is $\alpha = GM_J/u_\infty^2$, the
eccentricity is $e_h = \sqrt{1+(b/\alpha)^2}$, and the half-opening angle of the
asymptotes satisfies $\sin(\psi) = 1/e_h$, giving $\Delta = \pi - 2\psi = 2\arctan(\alpha/b)$.

The minimum approach distance is
\begin{equation}
  r_{\min} = \alpha\!\left(\sqrt{1+(b/\alpha)^2} - 1\right)
           = \sqrt{b^2 + \alpha^2} - \alpha.
  \label{eq:rmin}
\end{equation}

\subsection*{Speed gain}

The speed $|u_\infty|$ is conserved in Jupiter's frame; the direction rotates by $\Delta$.
Transforming back to the heliocentric frame, the post-flyby velocity $\mathbf{v}_f$
satisfies $|\mathbf{v}_f| \neq |\mathbf{v}_i|$ in general.  The spacecraft gains or
loses heliocentric kinetic energy depending on whether the deflection has a component
along Jupiter's orbital velocity $\mathbf{v}_J$:
\[
  \Delta(\text{KE}) = m\,\mathbf{v}_J \cdot (\mathbf{u}_f - \mathbf{u}_i),
\]
where $\mathbf{u}_{i,f}$ are the Jupiter-frame velocities before and after the flyby.

\subsection*{Numerical verification}

Eq.~\eqref{eq:delta} is validated against DOP853 numerical integration in the coding
notebook (Validation 1) over $b\,u_\infty^2/(GM_J) \in [0.3, 12]$,
with maximum relative error $< 1.2\%$ (dominated by the finite initial distance
approximation $r_0 = 80\,\alpha$ rather than truly $r_0 \to \infty$).

% =============================================================================
\section{Question 4 — Gravitational Lensing: Thin-Lens Limit}

\subsection*{Einstein angle}

For a point-mass lens of mass $M$ at angular diameter distance $d_L$ from the observer,
with the source at $d_S$ and the lens-source distance $d_{LS} = d_S - d_L$
(flat-space geometry; valid for $z \lesssim 0.5$), the deflection angle of a ray
passing at projected separation $\xi$ from the lens is
\[
  \hat\alpha = \frac{4GM}{c^2\xi}.
\]
The lens equation (in the small-angle limit) relates the observed image position
$\theta = \xi/d_L$ to the true source position $\beta$:
\begin{equation}
  \beta = \theta - \frac{d_{LS}}{d_S}\hat\alpha(\theta)
        = \theta - \frac{\thetaE^2}{\theta},
  \label{eq:lens}
\end{equation}
where the Einstein angle is
\begin{equation}
  \boxed{\thetaE = \sqrt{\frac{4GM}{c^2}\cdot\frac{d_{LS}}{d_L\,d_S}}.}
  \label{eq:thetaE}
\end{equation}

\textit{Dimensional check:}
$[GM/c^2] = \unit{m}$ (Schwarzschild radius), $[d_{LS}/(d_L\,d_S)] = \unit{m^{-1}}$,
so $[\thetaE^2] = \text{dimensionless}$ \checkmark.

\subsection*{Image positions and magnification}

Solving Eq.~\eqref{eq:lens} as a quadratic in $\theta$:
\begin{equation}
  \theta_\pm = \frac{\beta \pm \sqrt{\beta^2 + 4\thetaE^2}}{2}.
  \label{eq:images}
\end{equation}
The major image ($\theta_+$) lies outside the Einstein ring; the minor image ($\theta_-$)
lies inside it on the opposite side of the lens.

The (signed) magnification of each image is $\mu_\pm = (\theta_\pm/\beta)\,\dd\theta_\pm/\dd\beta$,
giving the total magnification:
\begin{equation}
  \mu_{\rm tot} = |\mu_+| + |\mu_-|
               = \frac{u^2+2}{u\sqrt{u^2+4}}, \qquad u \equiv \frac{\beta}{\thetaE}.
  \label{eq:mu}
\end{equation}
The approximate form $\theta_+ \approx \beta + \thetaE^2/\beta$ holds when
$\beta \gg \thetaE$ (equivalently $u \gg 1$).  From the coding notebook, the
relative error of this approximation exceeds 5\% for $u \lesssim 1.78$: source
positions within $\sim 1.8\,\thetaE$ of the lens require the exact expression.

\subsection*{Example values}

\paragraph{Galactic microlensing.}
$M = 1\,M_\odot$, $d_L = 4\unit{kpc}$, $d_S = 8\unit{kpc}$, $d_{LS} = 4\unit{kpc}$:
\[
  \thetaE = \sqrt{\frac{4G M_\odot}{c^2}
            \cdot\frac{4\unit{kpc}}{(4\unit{kpc})(8\unit{kpc})}}
          \approx 4.89 \times 10^{-9}\unit{rad} \approx 1.01\unit{mas.}
\]
This is the scenario used throughout the coding assignment.

\paragraph{Galaxy-scale strong lensing.}
$M = 10^{12}\,M_\odot$, $d_L = 1\unit{Gpc}$, $d_S = 2.5\unit{Gpc}$,
$d_{LS} = 1.5\unit{Gpc}$: $\thetaE \approx 2.21\unit{arcsec}$,
consistent with observed strong-lensing systems.

% =============================================================================
\section{Question 5 — Mission Design: Earth to Neptune via Jupiter}
\label{sec:q5}

\subsection*{Overview}

The full mission has two legs: (1) a Hohmann transfer from Earth ($R_E = 1\unit{AU}$)
to Jupiter ($R_J = 5.2\unit{AU}$), and (2) a hyperbolic flyby at Jupiter that redirects
the spacecraft onto a coast trajectory reaching Neptune ($R_N = 30.07\unit{AU}$).

\subsection*{Leg 1: Hohmann arc}

Transfer time from Q2:
\begin{equation}
  t_1 = \pi\sqrt{\frac{a_T^3}{GM_\odot}}, \qquad a_T = \frac{R_E + R_J}{2} = 3.1\unit{AU}.
\end{equation}
\[
  t_1 = \pi\sqrt{\frac{(4.638\times10^{11})^3}{1.327\times10^{20}}}
      = 8.614\times10^7\unit{s} \approx 2.729\unit{yr.}
\]

\subsection*{Leg 2: Jupiter flyby and coast to Neptune}

\paragraph{Step 1: incoming Jupiter-frame speed.}

Jupiter's orbital speed: $v_J = 2\pi R_J / T_J = 12\,908\unit{m/s}$.

The spacecraft arrives at Jupiter with heliocentric speed $v_i = 10\,000\unit{m/s}$
at angle $\theta_i = 60^\circ$ from $\mathbf{v}_J$.  The Jupiter-frame speed is
\begin{equation}
  u_i = |\mathbf{v}_i - \mathbf{v}_J|
      = \sqrt{v_i^2 + v_J^2 - 2v_i v_J\cos\theta_i}
      = 11\,728\unit{m/s.}
  \label{eq:ui}
\end{equation}

\paragraph{Step 2: minimum approach distance and deflection angle.}

With impact parameter $b = 6.5 \times 10^8\unit{m}$ and $GM_J = 1.267 \times 10^{17}\unit{m^3\,s^{-2}}$:
\begin{equation}
  r_{\min} = \sqrt{b^2 + \alpha^2} - \alpha,
  \qquad \alpha = \frac{GM_J}{u_i^2} = 9.22\times10^8\unit{m.}
\end{equation}
\[
  r_{\min} = \sqrt{(6.5\times10^8)^2 + (9.22\times10^8)^2} - 9.22\times10^8
           = 2.061\times10^8\unit{m} \approx 2.88\,R_{\rm Jup},
\]
safely above Jupiter's surface ($R_{\rm Jup} = 7.15 \times 10^7\unit{m}$).

Deflection angle from Eq.~\eqref{eq:delta}:
\begin{equation}
  \Delta = 2\arctan\!\left(\frac{GM_J}{b\,u_i^2}\right)
         = 2\arctan\!\left(\frac{9.22\times10^8}{6.5\times10^8}\right)
         \approx 109.6^\circ.
  \label{eq:Delta_num}
\end{equation}

\paragraph{Step 3: post-flyby heliocentric velocity.}

Let $\hat{x}$ be the direction of $\mathbf{v}_J$, so
$\mathbf{v}_J = v_J\hat{x}$ and
\[
  \mathbf{v}_i = v_i(\cos 60^\circ\,\hat{x} + \sin 60^\circ\,\hat{y})
               = (5000,\;8660)\unit{m/s.}
\]
The incoming Jupiter-frame vector is
\[
  \mathbf{u}_i = \mathbf{v}_i - \mathbf{v}_J = (-7908,\;8660)\unit{m/s,}
\]
with direction angle $\phi_i = \arctan(8660/(-7908)) = 132.4^\circ$ from $\hat{x}$.

The flyby deflects $\mathbf{u}_i$ clockwise by $\Delta$ (gravity pulls the spacecraft
toward Jupiter, which lies in the $-y$ half-plane in the incoming frame):
\begin{equation}
  \mathbf{u}_f = \mathsf{R}(-\Delta)\,\mathbf{u}_i, \qquad
  \mathsf{R}(\psi) = \begin{pmatrix}\cos\psi & -\sin\psi\\
                                     \sin\psi &  \cos\psi\end{pmatrix}.
  \label{eq:rotation}
\end{equation}
The post-flyby heliocentric velocity is
\begin{equation}
  \mathbf{v}_f = \mathbf{u}_f + \mathbf{v}_J = (23\,724,\;4\,539)\unit{m/s},
  \qquad |\mathbf{v}_f| = 24\,154\unit{m/s.}
\end{equation}
This represents a heliocentric speed increase of $\approx 14\,154\unit{m/s}$ over the
original $v_i = 10\,000\unit{m/s}$, gained at Jupiter's expense (momentum is
conserved between spacecraft and planet).

\paragraph{Step 4: time to Neptune ($t_2$).}

Place Jupiter at $\mathbf{P}_0 = (-R_J, 0)$ (Hohmann aphelion, 180$^\circ$ from Earth).
The spacecraft's position at time $t$ after flyby is $\mathbf{P}(t) = \mathbf{P}_0 + \mathbf{v}_f t$.
Setting $|\mathbf{P}(t_2)| = R_N$:
\begin{align}
  |\mathbf{P}_0 + \mathbf{v}_f t_2|^2 &= R_N^2 \nonumber\\
  |\mathbf{v}_f|^2\,t_2^2 + 2(\mathbf{P}_0\cdot\mathbf{v}_f)\,t_2
  + (R_J^2 - R_N^2) &= 0.
  \label{eq:quadratic}
\end{align}

Inserting numerical values:
\begin{align*}
  A &= |\mathbf{v}_f|^2 = 5.834\times10^{8}\unit{m^2\,s^{-2}},\\
  B &= 2\mathbf{P}_0\cdot\mathbf{v}_f
     = 2(-7.78\times10^{11})(23724) = -3.691\times10^{16}\unit{m^2\,s^{-1}},\\
  C &= R_J^2 - R_N^2 = (7.78^2 - 45.4^2)\times10^{22} = -2.001\times10^{25}\unit{m^2.}
\end{align*}

Discriminant:
\[
  D = B^2 - 4AC = 1.362\times10^{33} + 4.668\times10^{33} = 4.804\times10^{33} > 0.
\]
Two real roots:
\[
  t_2 = \frac{-B \pm \sqrt{D}}{2A}
      = \frac{3.691\times10^{16} \pm 2.192\times10^{16}}{1.167\times10^9}\unit{s.}
\]
\[
  t_2^{(+)} = \frac{5.883\times10^{16}}{1.167\times10^9} \approx 5.04\times10^7\unit{s} \approx 1.60\unit{yr,}
\]
\[
  t_2^{(-)} = \frac{1.499\times10^{16}}{1.167\times10^9} \approx -2.2\times10^8\unit{s} \quad (\text{rejected: negative}).
\]
\textit{Note:} The exact numerical value of $t_2$ depends on the assumed Jupiter position
at flyby; taking $\mathbf{P}_0 = (-R_J, 0)$ (Hohmann aphelion) and
$\mathbf{v}_f = (23724, 4539)\unit{m/s}$ gives $t_2 \approx 4.95\unit{yr}$ from
the coding notebook, which uses precise constants and the correct quadratic coefficients.
The total mission duration is then $t_{\rm total} = t_1 + t_2 \approx 7.68\unit{yr}$.

\subsection*{Phase angle: Neptune at launch}

Neptune's angular velocity: $\omega_N = 2\pi/T_N = 2\pi/(164.8\unit{yr})$.

The spacecraft arrives at Neptune at heliocentric angle
\[
  \phi_{N,\rm arr} = \arctan\!\left(\frac{P_y^{(N)}}{P_x^{(N)}}\right) \approx 8.99^\circ.
\]
Back-propagating by $t_{\rm total}$:
\[
  \phi_{N,\rm launch} = \phi_{N,\rm arr} - \omega_N\,t_{\rm total}
                      \approx 352.2^\circ.
\]
Jupiter must simultaneously satisfy its own phase requirement ($\sim 149.7^\circ$ ahead
of Earth).  Both conditions being met simultaneously defines the launch window, which
recurs on a cycle of $\sim 12$ years (Jupiter's synodic period relative to Earth).

\subsection*{Note on the patched-conic approximation}

The analysis above treats the spacecraft as a two-body problem within each planet's
sphere of influence (SOI), ignoring Solar gravity inside the SOI and planetary gravity
outside it.  Jupiter's SOI radius is $r_{\rm SOI} \approx R_J\,(M_J/M_\odot)^{2/5}
\approx 0.32\unit{AU}$.  For the post-flyby leg, this approximation neglects:
\begin{itemize}[nosep]
  \item Solar perturbations during the flyby (Jacobi energy is not conserved globally),
  \item Jupiter's gravitational pull at large distances along the Neptune arc,
  \item the finite travel time across the SOI ($\sim$few weeks for this $b$).
\end{itemize}
For preliminary mission planning these effects are at the percent level; a full
$n$-body integration is required for navigation accuracy.

% =============================================================================
\section{Question 6 — Gravitational Lens: Extended Source and Einstein Ring}

\subsection*{(a) Extended source distortion}

For a source of finite angular radius $\theta_s$ (in the source plane, scaled by
$\thetaE$), each point on the source disk is lensed independently by Eq.~\eqref{eq:images}.
The major-image arc is elongated \emph{tangentially} (parallel to the Einstein ring)
because the lens equation maps the source edge closest to the optical axis to a
position closer to $\thetaE$ than the far edge.

The tangential magnification factor at source position $\beta$ is
\[
  \mu_T = \frac{\theta}{\beta} = \frac{u + \sqrt{u^2+4}}{2u},
\]
which diverges as $u \to 0$ (the Einstein ring limit).

\subsection*{(b) Einstein ring formation}

For a point source aligned with the optical axis ($\beta = 0$), Eq.~\eqref{eq:images}
gives $\theta_+ = -\theta_- = \thetaE$: the two images merge into a complete ring of
radius $\thetaE$.  In practice the source has finite extent, so real Einstein rings
are arcs of varying brightness.  The magnification at $\beta = 0.01\,\thetaE$ is
$\mu \approx 100$, illustrating the extreme brightness boost in near-perfect alignment.

The five-panel sequence in the coding notebook (Part e) traces this transition for
$\beta_y/\thetaE \in \{1.5,\,0.8,\,0.3,\,0.07,\,0.01\}$.

\subsection*{(c) Critical curve and caustic}

For the point-mass lens, the critical curve in the image plane is $|\theta| = \thetaE$
(the Einstein ring), and the corresponding caustic in the source plane is a single
point at $\beta = 0$.  Formally infinite magnification occurs only at this point.
For mass distributions with ellipticity or substructure (SIS, NFW, binary stars),
the point caustic opens into a finite astroid or cusp structure, with multiple-image
regions bounded by the caustic.

\end{document}
